\section{Reproduction guide}
\label{app:reproduction}

This appendix records how to rebuild the paper and re-run every computational
artifact referenced above. It is written for a reader who wants to scrutinize
the results rather than take them on trust.

\subsection{Building this paper}

The manuscript is designed to build from a pinned, minimal TeX dependency set
without repository-specific absolute paths. From the repository root:

\begin{verbatim}
sh paper/build.sh
\end{verbatim}

The script changes into \texttt{paper/}, creates or reuses a local Guix profile
at \texttt{paper/.guix-profile} from \texttt{manifest.scm}, sources that
profile, and runs \texttt{latexmk} with \texttt{pdflatex} and \texttt{biber},
crossing references and the bibliography as needed. It writes
\texttt{paper/main.pdf} and the usual intermediate files, which are ignored by
version control.

As an alternative to the local-profile script, the manifest can be used with an
ephemeral Guix shell. From inside \texttt{paper/}, run:

\begin{verbatim}
guix shell -m manifest.scm -- \
  latexmk -pdf -interaction=nonstopmode -halt-on-error main.tex
\end{verbatim}

No absolute path appears in the sources.

\subsection{Repository verification artifacts}

The finite and computational claims correspond to artifacts in the repository:

\begin{itemize}[leftmargin=*]
  \item exact and binomial same-length witnesses:
        \texttt{scripts/} exact-rational verifiers and the Lean modules
        \texttt{AssemblyP1/FixedLengthExactCounterexample.lean} and
        \texttt{AssemblyP1/FixedLengthBinomialCounterexample.lean};
  \item Section~6.2 molecule witnesses:
        \texttt{scripts/verify\_se62\_bridging\_flow\_counterexample.py},
        \texttt{scripts/verify\_se62\_mb09\_bidirected\_graph.py},
        \texttt{scripts/verify\_samelength\_se62\_counterexample.py}, and the
        Lean modules \texttt{Section62BridgingCounterexample.lean} and
        \texttt{SameLengthSection62Counterexample.lean};
  \item same-length oriented rigidity:
        \texttt{scripts/verify\_oriented\_se62\_rigidity.py}, plus the
        kernel-checked abstract circulation core in
        \texttt{AssemblyP1/OrientedRigidity.lean} and circular-spectrum adapter
        in \texttt{AssemblyP1/WordLayer.lean}.  The repeat-theory bridge that
        discharges the primitive/periodic hypotheses remains separate, so these
        modules do not yet kernel-check the full $I_s$ rigidity theorem;
  \item equivalence coupling:
        \texttt{scripts/verify\_equivalence\_coupling.py}.
\end{itemize}

Every Python verifier is self-contained, uses exact \texttt{fractions.Fraction}
arithmetic, is deterministic, and exits non-zero on any failed assertion. The
Lean modules contain no \texttt{sorry}, \texttt{axiom}, or \texttt{admit};

\begin{verbatim}
lake build
\end{verbatim}

builds them under the pinned toolchain. Where a claim is kernel-checked, the
corresponding module is named in Appendix~\ref{app:ledger}; where a claim is
only computationally verified, it is labeled accordingly.

\subsection{Epistemic boundary of reproduction}

Reproduction verifies the \emph{finite mathematics}: spectra, likelihood ratios,
bridging certificates, and graph feasibility. It does not verify the
source-correspondence claims that identify a given finite objective or
candidate universe with the one the 2016 sentence intends, because those are
source-supported inferences rather than computations. The ledger in
Appendix~\ref{app:ledger} keeps the two apart.
